\documentclass[12pt]{article}
\usepackage[utf8]{inputenc}
\usepackage[T1]{fontenc}
\usepackage{hyperref}
\title{Lehali Language Resources for K-12 Curriculum Development}
\author{Compiled by Claude Sonnet 5 (claude-sonnet-5) for Andr\'as Kornai}
\date{}
\begin{document}
\maketitle

\section{Introduction}
Lehali (ISO 639-3 tql) is an Oceanic language of Vanuatu. It is spoken on the west coast of Ureparapara in the Banks Islands. It was formerly known as Teqel, and it is distinct from L\"oy\"op, spoken on the island's east coast. Sources give about 200 speakers (a 2010 figure), against about 400 in the survey brief. UNESCO lists it as Vulnerable. It belongs to the Torres--Banks group. Most of what is documented comes from the linguist Alexandre Fran\c{c}ois.

\section{Primary Educational Resources}
No school materials for Lehali turned up in the searches.

\section{Linguistic and Reference Materials}
\begin{itemize}
\item Wikipedia, Lehali language: \url{https://en.wikipedia.org/wiki/Lehali_language}
\item Wikipedia, Torres--Banks languages: \url{https://en.wikipedia.org/wiki/Torres%E2%80%93Banks_languages}
\item Glottolog, Lehali: \url{https://glottolog.org/resource/languoid/id/leha1243.bigmap.html}
\item Fran\c{c}ois, fieldwork and languages page: \url{https://marama.huma-num.fr/AF-field.htm}
\item Lehali entry in the TLOPO database: \url{https://tlopo.clld.org/languages/lehali}
\end{itemize}

\section{Religious and Community Resources}
None were found.

\section{Digital and Audio Resources}
None were found directly. The Fran\c{c}ois fieldwork page above may list recordings, but I did not check.

\section{Curriculum Development Context}
Lehali is documented mainly in linguistic work. I found no school materials. A curriculum would need new work with the community on Ureparapara, drawing on Fran\c{c}ois's documentation.
\end{document}